\documentclass[11pt,a4paper]{article}
\usepackage[margin=2.4cm]{geometry}
\usepackage{amsmath,amssymb}
\usepackage{booktabs}
\usepackage{microtype}
\usepackage{algorithm}
\usepackage{algpseudocode}
\usepackage[hidelinks]{hyperref}

\newcommand{\W}{\mathsf{W}}
\newcommand{\Lb}{\mathsf{L}}
\newcommand{\D}{\mathsf{D}}
\newcommand{\NN}{\mathbb{N}}
\DeclareMathOperator{\rank}{rank}
\DeclareMathOperator{\orb}{orb}
\DeclareMathOperator{\canon}{canon}
\DeclareMathOperator{\mover}{mover}
\DeclareMathOperator{\Steps}{Steps}
\DeclareMathOperator{\Caps}{Caps}
\DeclareMathOperator{\Preds}{Preds}

\title{How the Wolf and Lamb tables were built}
\author{}
\date{27 September 2026}

\begin{document}
\maketitle

\begin{abstract}
This note gives the exact algorithms behind the solution tables of Wolf and Lamb:
how a position is turned into a table index, how the game is cut into layers by
lamb count, the two breadth-first retrograde passes that assign every position its
value and exit distance under the 100-turn rule, the second solver that adds the
number of plies to the end of the game, the 2-bit packing used by the website,
and the independent verifier. Everything here is what the programs
\texttt{solver/solve.c}, \texttt{solver/dist.c}, \texttt{solver/pack.js} and
\texttt{solver/verify.js} actually do.
\end{abstract}

\section{States and the rules that matter}

Cells are numbered $0,\dots,24$ with cell $=5(r-1)+(c-1)$ for board coordinates
$(r,c)$, $r,c\in\{1,\dots,5\}$. A \emph{state} is $s=(w,\ell,\mu)$: $w$ the set
of the three wolf cells, $\ell$ the set of lamb cells, $\mu\in\{\W,\Lb\}$ the side
to move. Sets are 25-bit masks. The no-capture counter and the repetition history
are deliberately \emph{not} part of the state; Section~\ref{sec:value} explains
why the tables are nevertheless exact under both rules.

Rules used by the solver: a step moves a piece to an orthogonally adjacent empty
cell; a wolf may instead jump over one empty cell in a row or column onto a lamb,
which is removed; a side with no legal move loses; the wolves win when three or
fewer lambs remain; 100 plies without a capture is a draw; threefold repetition is
a draw.

\section{Indexing: from a state to a table entry}\label{sec:index}

\paragraph{Symmetry.} The eight symmetries of the square (the dihedral group
$D_4$) preserve every rule, so two states that differ by a symmetry have the same
value. Each symmetry is stored as a permutation of the 25 cells; to apply one to a
mask quickly, five lookup tables of 32 entries map each 5-bit row of the mask to
its image, and the images are OR-ed.

\paragraph{Wolf orbits.} The $\binom{25}{3}=2300$ wolf triples are enumerated
and ranked by the combinatorial number system,
\[
\rank_3(\{a<b<c\}) = \binom a1 + \binom b2 + \binom c3 \in \{0,\dots,2299\}.
\]
For each triple $w$ its \emph{canonical form} $\canon(w)=\min_{g\in D_4} g(w)$
(numerically smallest mask) is computed; distinct canonical forms are numbered
$o=0,\dots,318$ in order of first appearance. For every raw triple the solver
stores its orbit number $\orb(w)$ and the list of symmetries $g$ with
$g(w)=\canon(w)$ (its \emph{coset}; usually one element, up to eight when the
triple is itself symmetric).

\paragraph{Lamb rank.} Fix an orbit $o$ with canonical wolf mask $W_o$. The 22
non-wolf cells are numbered $0,\dots,21$ in increasing cell order; a $k$-subset
$\{c_0<c_1<\dots<c_{k-1}\}$ of them has colex rank
\[
\rank(\ell) = \sum_{i=0}^{k-1}\binom{c_i}{\,i+1\,} \in \Big\{0,\dots,\binom{22}{k}-1\Big\},
\]
which is a bijection onto that range and is monotone in the numeric value of the
compressed mask. Ranking costs one binomial lookup per lamb; unranking runs $i$
from $k-1$ down to $0$ and takes the largest $c$ with $\binom{c}{i+1}\le$ the
remaining rank.

\paragraph{Index.} For a state $(w,\ell,\mu)$ with $k=|\ell|$:
\begin{enumerate}
\item $o=\orb(w)$; among the coset symmetries $g$ choose the one minimising
$g(\ell)$ and put $\ell'=g(\ell)$ (so that a symmetric wolf triple still yields a
unique representative);
\item $\mathrm{idx} = \big(o\cdot\binom{22}{k}+\rank(\ell')\big)\cdot 2 + [\mu=\Lb]$.
\end{enumerate}
A layer with $k$ lambs therefore has $N_k = 319\cdot\binom{22}{k}\cdot 2$
entries. Two states get the same index if and only if they are symmetric images
of each other. Decoding an index back to a state (needed by the solver when it
processes a queue of indices) inverts the formula and unranks the lamb set.

\paragraph{Aliases.} When $W_o$ is fixed by a non-trivial symmetry, some lamb
masks $\ell'$ are not minimal under the stabiliser of $W_o$; their entries are
never the target of a lookup. The solver skips them during processing (Section
\ref{sec:pass}) and fills them at the end by copying the value of their canonical
twin, so that any index decodes to a valid entry. About 10\,\% of entries are
aliases.

\paragraph{Enumeration.} All lamb sets of a layer are visited in rank order by
enumerating the 22-bit masks with $k$ bits in increasing numeric order (Gosper's
hack), which coincides with colex order; the rank is then just a running counter.
A $2300$-entry table maps raw triples to orbits, and per orbit three
$256$-entry tables expand a compressed 22-bit mask to a 25-bit board mask.

\section{Layers}

Captures remove a lamb and nothing adds one, so the position graph is stratified
by $k$: within a layer only steps occur, and captures go from layer $k$ to layer
$k-1$. Layers with $k\le3$ are wolf wins by rule. The solver processes
$k=4,5,\dots,15$ in order and keeps in memory only the current layer and the
finished layer $k-1$.

\section{The value of a state under the rules}\label{sec:value}

For side $X$ and a state $s$ of layer $k$ define $t_X(c)$ for each legal
successor $c$ of $s$:
\[
t_X(c)=\begin{cases}
T_X(c) & c\in\Steps(s),\\
0 & c\in\Caps(s)\text{ and } V_{k-1}(c)=X,\\
\infty & c\in\Caps(s)\text{ and } V_{k-1}(c)\neq X,
\end{cases}
\]
and let $T_X$ be the least solution of
\begin{equation}\label{eq:T}
T_X(s)=\begin{cases}
0 & \text{no legal move, } \mover(s)\ne X,\\
\infty & \text{no legal move, } \mover(s)= X,\\
1+\min_c t_X(c) & \mover(s)=X,\\
1+\max_c t_X(c) & \mover(s)\ne X.
\end{cases}
\end{equation}
$T_X(s)$ is the number of plies within which $X$ can force, against any defence,
either a terminal win or a capture into a layer-$(k-1)$ state that is an $X$ win.
The value with a fresh no-capture counter is
\[
V_k(s)=\W\text{ if }T_\W(s)\le100,\quad \Lb\text{ if }T_\Lb(s)\le100,\quad \D\text{ otherwise}.
\]
$T_\W$ and $T_\Lb$ cannot both be finite. The two rules left out of the state are
handled by this definition: the counter resets at every capture, so ``within 100
plies inside the layer'' is exactly what the rule requires; and along a forced
line $T_X$ strictly decreases, so no position repeats, while a defender who has
$T_Y>100-c$ for the attacker can always keep it so and reach the clock (or a
repetition, also a draw). Hence $X$ wins from $s$ with $c$ plies on the counter
iff $T_X(s)\le 100-c$.

\section{The retrograde pass}\label{sec:pass}

Equation \eqref{eq:T} is solved by breadth-first search backwards from the
exits, once per side $X$. Values are assigned in increasing order of $T_X$, which
makes the first assignment to an $X$-to-move state its minimum and the assignment
to a defender state (made when its last successor has been valued) its maximum.
Algorithm~\ref{alg:T} is the exact procedure; \emph{distinct canonical} means
that successors or predecessors that map to the same table index are counted
once.

\begin{algorithm}[h]
\caption{One pass of \texttt{solve.c}: $T_X$ for layer $k$, side $X$}\label{alg:T}
\begin{algorithmic}[1]
\State $T[\cdot]\gets$ unset;\ $\mathit{cnt}[\cdot]\gets 0$;\ levels $Q_0,Q_1\gets\emptyset$
\ForAll{non-alias states $s$ of layer $k$} \Comment{init, parallel over orbits}
  \If{$\mover(s)=X$}
    \If{$s$ has a capture $c$ with $V_{k-1}(c)=X$} $T[s]\gets1$; push $s$ to $Q_1$ \EndIf
  \Else
    \If{$s$ has no legal move} $T[s]\gets0$; push $s$ to $Q_0$
    \ElsIf{$s$ has a capture $c$ with $V_{k-1}(c)\ne X$} $\mathit{cnt}[s]\gets\textsc{escape}$
    \ElsIf{$\Steps(s)=\emptyset$} $T[s]\gets1$; push $s$ to $Q_1$
    \Else\ $\mathit{cnt}[s]\gets|\{\text{distinct canonical } c\in\Steps(s)\}|$
    \EndIf
  \EndIf
\EndFor
\For{$d=0,1,2,\dots$ while $Q_d\cup Q_{d+1}\ne\emptyset$}
  \ForAll{$s\in Q_d$}
    \If{$d\le100$} $V_k[s]\gets X$ (error if already set by the other side) \EndIf
    \ForAll{distinct canonical $q\in\Preds(s)$} \Comment{$q$ reaches $s$ by a step}
      \If{$\mover(q)=X$ \textbf{and} $T[q]$ unset} $T[q]\gets d+1$; push $q$ to $Q_{d+1}$
      \ElsIf{$\mover(q)\ne X$ \textbf{and} $\mathit{cnt}[q]\ne\textsc{escape}$}
        $\mathit{cnt}[q]\gets\mathit{cnt}[q]-1$;\ \textbf{if} $\mathit{cnt}[q]=0$ \textbf{then} $T[q]\gets d+1$; push $q$ to $Q_{d+1}$
      \EndIf
    \EndFor
  \EndFor
\EndFor
\State after both sides: every non-alias state with $V_k$ unset is a draw; aliases copy their canonical twin
\end{algorithmic}
\end{algorithm}

\paragraph{Why the counters are right.} A defender state $q$ is valued when all
its step successors have been valued; the counter is initialised to the number
of distinct canonical successors and decremented once per processed successor
$c$ for which $q$ is a distinct canonical predecessor. These two counts agree: if
$q\to c_1$ is a step and $g$ a symmetry with $g(c_1)$ canonical, then
$g(q)\to g(c_1)$ is a step, so ``$c$ is a distinct canonical successor of $q$''
and ``$q$ is a distinct canonical predecessor of $c$'' are equivalent. Aliases
must be excluded from processing; otherwise a state and its mirror image would
both decrement the counters of their common predecessors.

\paragraph{Predecessors.} Within a layer the only backward move is undoing a
step: for each piece of the side that just moved, and each empty neighbour of
it, move the piece back. Captures are never undone inside a layer; they are the
exits, valued directly from layer $k-1$ during initialisation.

\paragraph{Cost.} Initialisation generates the moves of every state once
(parallel over wolf orbits, disjoint index ranges). Propagation touches each
valued state once and each in-layer edge once. Memory per layer: the value table
(2 bytes per entry), $T$ (2 bytes), counters (1 byte), the queues, and the
previous layer's values: about 3\,GB for the largest layer ($k=11$,
$450$ million entries). The whole solve took about 12 minutes on 12 threads.

\section{Plies to the end of the game}

\texttt{dist.c} reads the value tables and, for every won state with winner $X$,
computes
\[
D(s)=\begin{cases}0 & \text{no legal move},\\ 1+\min_c d(c) & \mover(s)=X,\\ 1+\max_c d(c) & \mover(s)\ne X,\end{cases}
\qquad
d(c)=\begin{cases} D(c) & \text{step},\\ D_{k-1}(c) & \text{capture},\\ 0 & \text{capture leaving three lambs},\end{cases}
\]
where captures are only followed into states that layer $k-1$ labels as $X$ wins.
Because capture successors carry arbitrary distances, plain level-by-level BFS
does not suffice; the pass uses a \emph{bucket queue} indexed by distance. An
$X$-to-move state is given a provisional distance from its best capture at
initialisation and is lowered whenever a step successor with a smaller distance
is processed, stale queue entries being skipped; a defender state's distance is
computed as $1+\max$ over all successors when its counter reaches zero (all
step successors valued by then, capture successors known). Processing buckets in
increasing order guarantees minimality for $X$-to-move states and exactness for
defender states.

The 100-turn rule is not modelled inside $D$. Instead a second level-BFS
computes $S(s)$, the capture-free stretch that a $D$-optimal strategy needs,
where the winner chooses among $D$-optimal moves the one minimising $S$ and only
lines whose captures land in positions that are themselves exact are counted
($S\le100$ in the layer below). If $S(s)\le100-c$ the $D$-optimal strategy is
legal under the clock and $D(s)$ is exact; otherwise $D(s)$ is a proven lower
bound and the tools display it with $\ge$. This affects $0.38\,\%$ of won states.

\section{Packing for the website}

\texttt{pack.js} reads each value table and writes 2 bits per entry
($1=\W$, $2=\Lb$, $3=\D$) in index order, entry $e$ at bit offset $2e$, into
files of at most 64\,MB (\texttt{layer\_k.p0.bin}, \texttt{layer\_k.p1.bin});
\texttt{index.json} records sizes and part boundaries. A lookup in the browser
recomputes the index of Section~\ref{sec:index} in JavaScript, fetches the 4\,KB
block containing byte $\lfloor e/4\rfloor$ with an HTTP range request, and
extracts the 2-bit field. No server is involved. All layers together are
651\,MB.

\section{Verification}

The solver is not trusted. \texttt{verify.js} re-implements the indexing
independently, generates moves with the bitboard engine (which is itself
checked move-for-move against the plain rules engine on 300\,000 positions), and
tests at every non-alias state a local condition on the labels. Call a child $c$
of $s$ a ``$Z$-win within $n$'' if it is reached by a capture into a layer-$(k-1)$
$Z$ win, or by a step to a state labelled $Z$ with $T\le n$.

\begin{center}\small
\begin{tabular}{@{}lp{10.5cm}@{}}
\toprule
label of $s$ & required property \\
\midrule
no legal move & label $=$ opponent, $T=0$ \\
mover $M$ wins, $T$ & some child is an $M$-win within $T-1$, none within $T-2$, $1\le T\le100$ \\
opponent $Y$ wins, $T$ & every child is a $Y$-win within $T-1$, some not within $T-2$, $1\le T\le100$ \\
draw & no child is an $M$-win within 99, some child is not a $Y$-win within 99 \\
\bottomrule
\end{tabular}
\end{center}

By induction on $T$ these conditions force the labels to equal the true values,
given a correct layer $k-1$ and the base rule for three lambs. Aliases are
checked to equal their twins, and the indexing is self-tested on random positions
(all eight symmetric images map to one index, and the index decodes to an
equivalent state). \texttt{verify\_dist.js} checks $D$ and $S$ by the analogous
local conditions. A separate spot check compares random entries with a
brute-force search over the plain rules. All checks reported zero errors.

\section{File formats}

\begin{center}\small
\begin{tabular}{@{}llp{8.2cm}@{}}
\toprule
file & entry & meaning \\
\midrule
\texttt{solver/tables/layer\_k.bin} & uint16 & bits 15--14: value (1 $\W$, 2 $\Lb$, 3 $\D$); bits 13--0: $T$ of the winner \\
\texttt{solver/tables/dist\_k.bin} & uint16 & $D$ (0xFFFF for draws) \\
\texttt{solver/tables/seg\_k.bin} & uint16 & $S$ (0xFFFF for draws, or when no exact $D$-optimal line exists) \\
\texttt{tables/layer\_k.pN.bin} & 2 bits & value only, for the website \\
\bottomrule
\end{tabular}
\end{center}

Entry order in every file is the index of Section~\ref{sec:index}; all integers
are little-endian.

\end{document}
